% Options for packages loaded elsewhere
\PassOptionsToPackage{unicode}{hyperref}
\PassOptionsToPackage{hyphens}{url}
\PassOptionsToPackage{dvipsnames,svgnames,x11names}{xcolor}
\documentclass[
  10pt,
]{article}
\usepackage{xcolor}
\usepackage[margin=2.4cm]{geometry}
\usepackage{amsmath,amssymb}
\setcounter{secnumdepth}{-\maxdimen} % remove section numbering
\usepackage{iftex}
\ifPDFTeX
  \usepackage[T1]{fontenc}
  \usepackage[utf8]{inputenc}
  \usepackage{textcomp} % provide euro and other symbols
\else % if luatex or xetex
  \usepackage{unicode-math} % this also loads fontspec
  \defaultfontfeatures{Scale=MatchLowercase}
  \defaultfontfeatures[\rmfamily]{Ligatures=TeX,Scale=1}
\fi
\usepackage{lmodern}
\ifPDFTeX\else
  % xetex/luatex font selection
  \setmainfont[]{Libertinus Serif}
  \setmonofont[Scale=0.85]{Latin Modern Mono}
  \setmathfont[]{Latin Modern Math}
\fi
% Use upquote if available, for straight quotes in verbatim environments
\IfFileExists{upquote.sty}{\usepackage{upquote}}{}
\IfFileExists{microtype.sty}{% use microtype if available
  \usepackage[]{microtype}
  \UseMicrotypeSet[protrusion]{basicmath} % disable protrusion for tt fonts
}{}
\makeatletter
\@ifundefined{KOMAClassName}{% if non-KOMA class
  \IfFileExists{parskip.sty}{%
    \usepackage{parskip}
  }{% else
    \setlength{\parindent}{0pt}
    \setlength{\parskip}{6pt plus 2pt minus 1pt}}
}{% if KOMA class
  \KOMAoptions{parskip=half}}
\makeatother
\usepackage{longtable,booktabs,array}
\newcounter{none} % for unnumbered tables
\usepackage{calc} % for calculating minipage widths
% Correct order of tables after \paragraph or \subparagraph
\usepackage{etoolbox}
\makeatletter
\patchcmd\longtable{\par}{\if@noskipsec\mbox{}\fi\par}{}{}
\makeatother
% Allow footnotes in longtable head/foot
\IfFileExists{footnotehyper.sty}{\usepackage{footnotehyper}}{\usepackage{footnote}}
\makesavenoteenv{longtable}
\setlength{\emergencystretch}{3em} % prevent overfull lines
\providecommand{\tightlist}{%
  \setlength{\itemsep}{0pt}\setlength{\parskip}{0pt}}

\providecommand{\E}{\mathbb{E}}
\providecommand{\N}{\mathbb{N}}
\providecommand{\Z}{\mathbb{Z}}
\providecommand{\R}{\mathbb{R}}
\providecommand{\eps}{\varepsilon}
\providecommand{\ceil}[1]{\left\lceil #1\right\rceil}
\providecommand{\floor}[1]{\left\lfloor #1\right\rfloor}
\AtBeginDocument{\let\setminus\smallsetminus}
\usepackage[most]{tcolorbox}
\newtcolorbox{warning}{colback=red!5,colframe=red!65!black,boxrule=0.8pt,arc=2pt,left=6pt,right=6pt,top=5pt,bottom=5pt}
\usepackage{etoolbox}
\AtBeginEnvironment{longtable}{\small}
\setlength{\emergencystretch}{3em}
\usepackage{fancyhdr}
\pagestyle{fancy}
\fancyhf{}
\fancyhead[L]{\small\bfseries UNREFEREED GPT-6 ASTRA OUTPUT}
\fancyhead[R]{\small\thepage}
\renewcommand{\headrulewidth}{0.4pt}
\usepackage{bookmark}
\IfFileExists{xurl.sty}{\usepackage{xurl}}{} % add URL line breaks if available
\urlstyle{same}
\hypersetup{
  pdftitle={Conjecture 1.2},
  pdfauthor={Writeup: gpt-6-astra (single model pass)},
  colorlinks=true,
  linkcolor={blue!50!black},
  filecolor={Maroon},
  citecolor={Blue},
  urlcolor={blue!50!black},
  pdfcreator={LaTeX via pandoc}}

\title{Conjecture 1.2}
\usepackage{etoolbox}
\makeatletter
\providecommand{\subtitle}[1]{% add subtitle to \maketitle
  \apptocmd{\@title}{\par {\large #1 \par}}{}{}
}
\makeatother
\subtitle{Unrefereed candidate disproof by GPT-6 Astra}
\author{Writeup: \texttt{gpt-6-astra} (single model pass)}
\date{Catalog id \texttt{2509.09031\_\_00} --- generated 2026-09-17}

\begin{document}
\maketitle

\begin{warning}

\textbf{UNREFEREED MODEL OUTPUT.} This document was selected solely
because GPT-6 Astra labelled its own result
\texttt{would\_publish:\ true} and returned \texttt{proved} or
\texttt{disproved}. It has not passed the adversarial LLM referee used
for the earlier Sol campaign, has not been checked by a human
mathematician, and has not been checked for novelty. Treat every
mathematical and bibliographic claim below as unverified.

\end{warning}

\section{Summary and provenance}\label{summary-and-provenance}

{\def\LTcaptype{none} % do not increment counter
\begin{longtable}[]{@{}
  >{\raggedright\arraybackslash}p{(\linewidth - 2\tabcolsep) * \real{0.5000}}
  >{\raggedright\arraybackslash}p{(\linewidth - 2\tabcolsep) * \real{0.5000}}@{}}
\toprule\noalign{}
\endhead
\bottomrule\noalign{}
\endlastfoot
Catalog id & \texttt{2509.09031\_\_00} \\
Catalog entry &
\href{https://graph-theory-ai.github.io/graph-conjectures/arxiv/2509.09031__00/}{Conjecture
1.2} \\
Source corpus & arxiv \\
Campaign leg & selected second attempts (\texttt{attacks\_retry}) \\
Source paper / entry & Asymptotic structure. II. Path-width and additive
quasi-isometry \\
Model verdict & \textbf{disproved} (confidence: high) \\
Model's one-line claim & A finite counterexample is obtained from a
contraction- and subdivision-closed class whose branch vertices are
pairwise joined by internally degree-two paths. \\
Model & \texttt{gpt-6-astra}, reasoning effort \texttt{max},
\texttt{mode=pro}, flex service tier \\
Original artifact &
\texttt{attacks\_retry/2509.09031\_\_00/output.md} \\
Independent review & \textbf{None} \\
Model's caveats & Uses the standard coarsely onto definition of
quasi-isometry; the class is not required to be closed under edge
deletion. \\
\end{longtable}
}

\section{Problem statement}\label{problem-statement}

\emph{No extracted statement was stored in the artifact metadata.}

\section{Astra writeup}\label{astra-writeup}

\emph{The text below is the model output verbatim, apart from moving its
machine-readable verdict block into the summary above and shifting
Markdown heading levels for this document. Mechanical TeX defects and
equation tags were normalized where needed for compilation.}

\subsection{A finite counterexample}\label{a-finite-counterexample}

All graphs below are finite, connected, and simple. Edge contraction
includes the usual suppression of loops and parallel edges.

An \((L,C)\)-quasi-isometry \(f:G\to H\) satisfies \[
L^{-1}d_G(x,x')-C\le d_H(f(x),f(x'))\le Ld_G(x,x')+C
\] and every vertex of \(H\) is within distance \(C\) of \(f(V(G))\).

We construct a fixed class \(\mathcal C\), closed under edge contraction
and subdivision, and graphs \(G_g,H_g^+\) such that:

\begin{itemize}
\tightlist
\item
  \(H_g^+\in\mathcal C\);
\item
  there is an onto \((2,1)\)-quasi-isometry \(G_g\to H_g^+\);
\item
  every \((1,A)\)-quasi-isometry from \(G_g\) to any member of
  \(\mathcal C\) satisfies \[
  A+1>\frac{g}{10000}.
  \]
\end{itemize}

Thus the proposed uniform additive constant does not exist.

The red/blue turn-cost mechanism from the supplied attempt is useful,
but its arbitrary-map/arbitrary-target gap must be closed. We do this by
adding long complete-graph corridors and proving a rigidity lemma. The
infinite-graph argument is not used.

\begin{center}\rule{0.5\linewidth}{0.5pt}\end{center}

\subsubsection{1. The contraction- and subdivision-closed
class}\label{the-contraction--and-subdivision-closed-class}

For a graph \(J\), put \[
B(J)=\{v\in V(J):\deg_J(v)\ge 3\}.
\] A \textbf{thread} between distinct vertices \(u,v\in B(J)\) is a
\(u\)-\(v\) path all of whose internal vertices have degree two in
\(J\).

Define \[
\mathcal C=
\left\{
J:\text{every two distinct vertices of }B(J)
\text{ are joined by a thread}
\right\}.
\qquad\text{(1.1)}
\] In particular, every subdivision of a complete graph belongs to
\(\mathcal C\).

\paragraph{Lemma 1.1}\label{lemma-1.1}

The class \(\mathcal C\) is closed under edge subdivision and edge
contraction.

\subparagraph{Proof}\label{proof}

Subdivision is immediate: the old branch vertices are unchanged, and
threads remain threads.

For contraction, let \(J'\) be obtained from \(J\in\mathcal C\) by
contracting \(xy\). Every branch vertex of \(J'\) is the image of a
branch vertex of \(J\). Indeed, degrees of uncontracted vertices cannot
increase, and contracting two vertices of degree at most two produces a
vertex of degree at most two.

Take distinct \(u',v'\in B(J')\), and choose branch vertices
\(u,v\in B(J)\) mapping to them. Let \(P\) be a thread between \(u\) and
\(v\) in \(J\).

The image of \(P\) contains a \(u'\)-\(v'\) path whose internal vertices
have degree at most two in \(J'\). To see the only potentially
problematic point, suppose the contracted vertex occurs internally on
this image. If one endpoint of \(xy\) is internal on \(P\), its two
neighbours already lie on \(P\). Hence the other endpoint of \(xy\) also
lies on \(P\); unless the contracted vertex is an endpoint of the image,
both contracted vertices have degree two. Their contraction therefore
has degree at most two.

After choosing a simple path in the image, its internal vertices have
degree exactly two. Thus \(u',v'\) are joined by a thread. \(\square\)

\begin{center}\rule{0.5\linewidth}{0.5pt}\end{center}

\subsubsection{2. Metric preliminaries}\label{metric-preliminaries}

We use the geometric realization of a graph, giving every edge length
one. Distances between original vertices are unchanged.

\paragraph{Lemma 2.1 --- Continuous coarse
inverses}\label{lemma-2.1-continuous-coarse-inverses}

If there is a \((1,A)\)-quasi-isometry between graphs \(G\) and \(K\),
then their geometric realizations \(X,Y\) admit continuous maps \[
F:X\to Y,\qquad Q:Y\to X
\] such that, with \[
D=10(A+1),
\qquad\text{(2.1)}
\] we have \begin{equation*}
\begin{aligned}
|d_Y(Fx,Fx')-d_X(x,x')|&\le D,\\
|d_X(Qy,Qy')-d_Y(y,y')|&\le D,\\
d_Y(FQy,y)&\le D.
\end{aligned}
\qquad\text{(2.2)}
\end{equation*}

\subparagraph{Proof}\label{proof-1}

Extend the given vertex map over each edge along a geodesic. The
resulting continuous map \(F\) has distortion at most \(3A+4\): each
point can be compared with an endpoint of its edge, and each image edge
has length at most \(A+1\).

For every vertex \(w\) of \(K\), choose a vertex \(q(w)\) of \(G\) with
\[
d_K(f(q(w)),w)\le A.
\] The vertex map \(q\) has distortion at most \(3A\). Extend it over
edges along geodesics. The resulting \(Q\) has distortion at most
\(9A+4\).

Finally, comparing \(y\) with an endpoint \(w\) of its edge gives \[
d_Y(FQy,y)\le 7A+6.
\] All three bounds are at most \(10(A+1)\). \(\square\)

We also use the following elementary observation.

\paragraph{Lemma 2.2 --- Detecting a branch
vertex}\label{lemma-2.2-detecting-a-branch-vertex}

Let \(Z\) be a metric graph. Suppose \(z_1,z_2,z_3\) lie within distance
\(r\) of \(z\), and \[
d_Z(z_i,z_j)>r\qquad(i\ne j).
\] Then \(B(Z)\) meets the closed radius-\(r\) ball about \(z\).

\subparagraph{Proof}\label{proof-2}

Otherwise, the union of shortest \(z\)-\(z_i\) paths has no branch
vertex. It lies in a path, or forms a cycle of length at most \(2r\). In
the path case, two of the three points lie on the same side of \(z\), at
distance at most \(r\) from each other. In the cycle case, the diameter
is at most \(r\). Both contradict the hypothesis. \(\square\)

\begin{center}\rule{0.5\linewidth}{0.5pt}\end{center}

\subsubsection{3. Rigidity of a complete system of long
corridors}\label{rigidity-of-a-complete-system-of-long-corridors}

Here is the step that rules out changing either the quasi-isometry or
the target.

\paragraph{Lemma 3.1 --- Long-corridor
rigidity}\label{lemma-3.1-long-corridor-rigidity}

Let \(n\ge4\), \(D\ge1\), and \(S\ge1000D\).

Construct a metric graph \(X\) as follows:

\begin{itemize}
\tightlist
\item
  replace each vertex \(v\) of \(K_n\) by a tree \(T_v\) of diameter at
  most one;
\item
  choose \(r_v\in T_v\);
\item
  for each pair \(u,v\), join a point of \(T_u\) to a point of \(T_v\)
  by one corridor \(I_{uv}\), of length \[
  L_{uv}\ge S;
  \]
\item
  these corridors have disjoint interiors, and there are no other edges.
\end{itemize}

Let \(Y\) be the geometric realization of a graph in \(\mathcal C\).
Suppose continuous maps \(F:X\to Y\), \(Q:Y\to X\) satisfy (2.2).

Then:

\begin{enumerate}
\def\labelenumi{\arabic{enumi}.}
\tightlist
\item
  \(Y\) has exactly \(n\) branch vertices, labelled \(h_v\);
\item
  for each pair \(u,v\), there is exactly one thread between
  \(h_u,h_v\);
\item
  writing \(\lambda_{uv}\) for its length, the weighted complete-graph
  metric \(d_\lambda\) satisfies \[
  |d_\lambda(u,v)-d_X(r_u,r_v)|\le30D;
  \qquad\text{(3.1)}
  \]
\item
  every such thread satisfies \[
  \lambda_{uv}\ge \frac{L_{uv}}2-D.
  \qquad\text{(3.2)}
  \]
\end{enumerate}

\subparagraph{Proof}\label{proof-3}

\paragraph{\texorpdfstring{Step 1: Find a branch vertex near each
\(F(r_v)\)}{Step 1: Find a branch vertex near each F(r\_v)}}\label{step-1-find-a-branch-vertex-near-each-fr_v}

Choose three different corridors incident with \(T_v\), and on each
choose a point at corridor-distance \(10D\) from \(T_v\). Each point is
within \(10D+1\) of \(r_v\), and each pair is at distance at least
\(20D\).

Their images under \(F\) lie within \(12D\) of \(F(r_v)\), and have
pairwise distances at least \(19D\). Lemma 2.2 therefore gives \[
h_v\in B(Y),\qquad d_Y(h_v,F(r_v))\le12D.
\qquad\text{(3.3)}
\]

For distinct \(u,v\), we have \(d_X(r_u,r_v)\ge S\), so \[
d_Y(h_u,h_v)\ge S-25D.
\qquad\text{(3.4)}
\] In particular, these \(n\) branch vertices are distinct.

\paragraph{\texorpdfstring{Step 2: Every branch vertex of \(Y\) belongs
to one of these
clusters}{Step 2: Every branch vertex of Y belongs to one of these clusters}}\label{step-2-every-branch-vertex-of-y-belongs-to-one-of-these-clusters}

Let \(h\in B(Y)\). By (3.4), at most one of the chosen \(h_v\) lies
within distance \(S/3\) of \(h\). Since \(n\ge4\), choose three others,
each at distance at least \(S/3\).

The defining property of \(\mathcal C\) supplies threads from \(h\) to
these three vertices. The threads have disjoint interiors, and each has
length at least \(S/3\). Points at thread-distance \(10D\) from \(h\)
therefore have distance \(10D\) from \(h\) and pairwise distance
\(20D\).

Their images under \(Q\), together with Lemma 2.2, show that \(Q(h)\)
lies within \(11D\) of a branch vertex of \(X\). All branch vertices of
\(X\) lie in the trees \(T_v\). Consequently, \[
d_X(Q(h),r_v)\le12D
\qquad\text{(3.5)}
\] for some \(v\). Using (2.2), \[
d_Y(h,F(r_v))\le14D.
\qquad\text{(3.6)}
\] This \(v\) is unique because distinct \(F(r_v)\) are at distance at
least \(S-D>28D\).

Call it the cluster of \(h\). The initially chosen \(h_v\) belongs to
cluster \(v\).

\paragraph{Step 3: A thread between two clusters must follow the
corresponding
corridor}\label{step-3-a-thread-between-two-clusters-must-follow-the-corresponding-corridor}

Consider any thread \(P\) between branch vertices \(h,h'\) in distinct
clusters \(v,w\). Its length \(\lambda\) satisfies \[
\lambda\ge d_Y(h,h')\ge S-29D>100D.
\] Remove the first and last \(50D\) of \(P\), leaving a closed interval
\(P^0\).

Every point \(y\in P^0\) has distance at least \(50D\) from \(B(Y)\),
since the interior of a thread can be left only through its endpoints.

If \(d_X(Qy,r_z)\le20D\), then (2.2) and (3.3) give \[
d_Y(y,h_z)
 \le D+(20D+D)+12D
 =34D,
\] a contradiction. Thus \[
Q(P^0)\cap\bigcup_z\overline B_X(r_z,20D)=\varnothing.
\qquad\text{(3.7)}
\]

The components of the complement in (3.7) are precisely trimmed
interiors of the individual corridors of \(X\). Since \(Q\) is
continuous, \(Q(P^0)\) lies in one such component.

Let \(y_0,y_1\) be the two endpoints of \(P^0\). By (3.5), \[
d_X(Qy_0,r_v)\le 50D+D+12D=63D,
\] and similarly \(d_X(Qy_1,r_w)\le63D\). Since \(63D<S\), the corridor
containing \(Q(P^0)\) must be incident with both \(v\) and \(w\). It is
therefore \(I_{vw}\).

Moreover, measured from the respective ends of \(I_{vw}\), the points
\(Qy_0,Qy_1\) lie within \(63D\). By continuity, \(Q(P^0)\) covers the
whole central interval \[
[63D,L_{vw}-63D]
\qquad\text{(3.8)}
\] of that corridor.

\paragraph{Step 4: No two threads can join the same two
clusters}\label{step-4-no-two-threads-can-join-the-same-two-clusters}

If two distinct threads joined clusters \(v,w\), (3.8) would give
interior points \(y,y'\), one on each trimmed thread, mapping under
\(Q\) to the midpoint of \(I_{vw}\).

The distortion bound for \(Q\) would imply \[
d_Y(y,y')\le D.
\] But distinct threads have disjoint interiors, and leaving the first
trimmed thread requires travelling at least \(50D\). This is impossible.

Every cluster is nonempty. If some cluster contained two branch
vertices, threads from both of them to a branch vertex in any other
cluster would contradict the preceding paragraph. Thus every cluster
contains exactly one branch vertex.

The defining property of \(\mathcal C\) now gives exactly one thread
between each pair \(h_u,h_v\).

\paragraph{Step 5: The metric and thread-length
conclusions}\label{step-5-the-metric-and-thread-length-conclusions}

Any shortest path between branch vertices traverses a sequence of the
threads just identified. Other components---pendant paths or cycles
attached at a single branch vertex---cannot shorten such a path.
Therefore their metric is exactly the weighted \(K_n\) metric with edge
lengths \(\lambda_{uv}\).

Equation (3.3) and the distortion of \(F\) give \[
|d_Y(h_u,h_v)-d_X(r_u,r_v)|\le25D,
\] which implies (3.1).

Finally, (3.8) includes the quarter and three-quarter points of
\(I_{uv}\). Their distance in \(X\) is \(L_{uv}/2\): the direct interval
has that length, and a route through the two ends cannot be shorter.
Choosing preimages on the thread and using the distortion of \(Q\)
yields \[
\lambda_{uv}\ge L_{uv}/2-D.
\] This proves (3.2). \(\square\)

\begin{center}\rule{0.5\linewidth}{0.5pt}\end{center}

\subsubsection{4. High-girth graphs with two
colours}\label{high-girth-graphs-with-two-colours}

For every integer \(g\ge5\), we need a finite connected simple graph
\(H_g\) of girth at least \(g\), with \[
E(H_g)=E_R\sqcup E_B,
\] such that every vertex has two incident red edges and two incident
blue edges.

Here is a self-contained construction.

Let \(W_g\) be the finite set of all nonempty freely reduced words of
length less than \(g\) in \[
a,a^{-1},b,b^{-1}.
\] For a word \(w\) of length \(k\), take the set \(\{0,\ldots,k\}\),
and prescribe partial permutations \(a_w,b_w\) so that following the
letters of \(w\) sends \[
0\to1\to\cdots\to k.
\] For an inverse letter, prescribe the corresponding reversed
permutation arrow. These partial permutations are well-defined and
injective: a conflict at an intermediate index would require adjacent
inverse letters. Complete each partial permutation to a permutation, for
example by matching the remaining domain and range elements in
increasing order.

In the product of these finite symmetric groups, let \(a,b\) be the
tuples of the constructed permutations, and let \(Q_g\) be the subgroup
they generate. No nonempty freely reduced word of length less than \(g\)
is trivial in \(Q_g\), because its own coordinate sends \(0\) to \(k\).

The undirected Cayley graph of \(Q_g\) with generators
\(a^{\pm1},b^{\pm1}\) is therefore simple, connected, 4-regular, and has
girth at least \(g\). Colour \(a\)-edges red and \(b\)-edges blue.

Each monochromatic subgraph is a union of cycles, all of length at least
\(g\).

\begin{center}\rule{0.5\linewidth}{0.5pt}\end{center}

\subsubsection{5. The graphs and their uniform
quasi-isometries}\label{the-graphs-and-their-uniform-quasi-isometries}

Fix an integer \(g\ge10000\), and write \[
H=H_g,\qquad n=|V(H)|,\qquad S=g,\qquad M=100g^2.
\qquad\text{(5.1)}
\]

Construct \(H_g^+\) from \(K_n\), on vertex set \(V(H)\), by replacing:

\begin{itemize}
\tightlist
\item
  every edge of \(H\) by a path of length \(S\);
\item
  every edge not in \(H\) by a path of length \(M\).
\end{itemize}

Thus \(H_g^+\) is a subdivision of \(K_n\), so \[
H_g^+\in\mathcal C.
\]

Construct \(G_g\) by splitting each original vertex \(v\) into two
vertices \(v_R,v_B\), joined by a unit-length \textbf{switch edge}.
Attach:

\begin{itemize}
\tightlist
\item
  the red \(H\)-corridors to the corresponding \(R\)-vertices;
\item
  the blue \(H\)-corridors to the corresponding \(B\)-vertices;
\item
  every length-\(M\) corridor to the corresponding \(R\)-vertices at
  both ends.
\end{itemize}

Define \[
p:G_g\to H_g^+
\] by collapsing each switch edge and mapping all corridor vertices to
their corresponding vertices.

This map is onto. Every \(G_g\)-path projects to an \(H_g^+\)-walk of no
greater length, so \[
d_{H_g^+}(p(x),p(y))\le d_{G_g}(x,y).
\] Conversely, a shortest \(H_g^+\)-path of length \(k\) lifts using at
most \(k+1\) switches. Hence \[
d_{G_g}(x,y)\le2d_{H_g^+}(p(x),p(y))+1.
\] Therefore \(p\) is an onto \((2,1)\)-quasi-isometry.

\begin{center}\rule{0.5\linewidth}{0.5pt}\end{center}

\subsubsection{6. No uniformly additive approximation
exists}\label{no-uniformly-additive-approximation-exists}

Suppose there is a \((1,A)\)-quasi-isometry \[
G_g\to K,\qquad K\in\mathcal C,
\] where \[
A+1\le \frac{g}{10000}.
\qquad\text{(6.1)}
\] Put \[
D=10(A+1),\qquad E=30D,\qquad
s=\left\lfloor\frac g{10}\right\rfloor.
\qquad\text{(6.2)}
\] Then \(S=g\ge1000D\), so Lemmas 2.1 and 3.1 apply to \(G_g\), with
the switch edges as the trees \(T_v\) and \(r_v=v_R\).

We obtain positive lengths \(\lambda_{uv}\) on the edges of \(K_n\) such
that \[
|d_\lambda(u,v)-d_{G_g}(u_R,v_R)|\le E
\qquad\text{(6.3)}
\] for every \(u,v\), and \[
\lambda_{uv}\ge M/2-D
\qquad\text{if }uv\notin E(H).
\qquad\text{(6.4)}
\]

The following inequalities follow from \(g\ge10000\) and
\(D\le g/1000\): \begin{equation*}
\begin{gathered}
s-3>2E,\qquad Ss>2+2E,\qquad S>s+1,\\
M/2-D>Ss+2+E,\qquad 2s<g/2,\qquad g-s\ge2s.
\qquad\text{(6.5)}
\end{gathered}
\end{equation*}

We now show that (6.3)--(6.4) are impossible.

\paragraph{\texorpdfstring{6.1. Lower bounds on short
\(H\)-paths}{6.1. Lower bounds on short H-paths}}\label{lower-bounds-on-short-h-paths}

If \(P\) is an \(H\)-path of length \(t<g/2\), then \(P\) is geodesic in
\(H\). Let its endpoints be \(u,v\).

Any \(G_g\)-path from \(u_R\) to \(v_R\) that uses a long corridor has
length at least \(M>St\). A path using only \(H\)-corridors has length
at least \(St\). Thus \[
d_{G_g}(u_R,v_R)\ge St.
\] Consequently, \[
\lambda(P)\ge d_\lambda(u,v)\ge St-E.
\qquad\text{(6.6)}
\]

\paragraph{6.2. Upper bounds on monochromatic
paths}\label{upper-bounds-on-monochromatic-paths}

Let \(P\) be a monochromatic \(H\)-path of length \(s\), with endpoints
\(u,v\).

If \(P\) is red, it lifts from \(u_R\) to \(v_R\) with length \(Ss\). If
it is blue, two endpoint switches suffice. Hence \[
d_{G_g}(u_R,v_R)\le Ss+2,
\] and so \[
d_\lambda(u,v)\le Ss+2+E.
\qquad\text{(6.7)}
\]

By (6.4)--(6.5), a \(\lambda\)-shortest path \(Q\) between \(u,v\)
cannot use an edge outside \(H\). Choose \(Q\) simple.

If \(Q\ne P\), their union contains a cycle, giving \[
|Q|+s\ge g.
\] In particular, \(Q\) contains a \(2s\)-edge subpath. Applying (6.6)
to that subpath gives \[
\lambda(Q)\ge2Ss-E>Ss+2+E,
\] contrary to (6.7). Thus \(P\) itself is shortest, and \[
\lambda(P)\le Ss+2+E.
\qquad\text{(6.8)}
\]

\paragraph{6.3. Average over the monochromatic
cycles}\label{average-over-the-monochromatic-cycles}

Apply (6.8) to every cyclic window of \(s\) edges on every monochromatic
cycle. Each edge of a cycle occurs in exactly \(s\) windows. Summing
over both colours gives \[
\sum_{e\in E(H)}\lambda_e
\le
|E(H)|\left(S+\frac{2+E}{s}\right).
\qquad\text{(6.9)}
\]

Now choose a uniformly random directed edge of \(H\), and continue for
\(s\) edges by alternating colours, choosing uniformly between the two
available edges of the required colour.

At every position the directed edge is uniformly distributed: every
directed edge has exactly two possible predecessor directed edges of the
opposite colour, each transitioning to it with probability \(1/2\).

Thus (6.9) implies that the expected \(\lambda\)-length of this walk is
at most \[
Ss+2+E.
\] There is therefore an alternating walk \[
P=v_0v_1\cdots v_s
\] with \[
\lambda(P)\le Ss+2+E.
\qquad\text{(6.10)}
\] It is nonbacktracking. Since \(s<g\), it is a path; since \(2s<g\),
it is the unique \(H\)-geodesic between its endpoints.

\paragraph{6.4. Alternating paths require accumulated switch
costs}\label{alternating-paths-require-accumulated-switch-costs}

Following \(P\) in \(G_g\), from \((v_0)_R\) to \((v_s)_R\), costs at
most \[
Ss+s+1.
\] A shortest path cannot use a long corridor, since \(M\) exceeds this
quantity.

Nor can it use at least \(s+1\) short corridors, since \[
S(s+1)>Ss+s+1.
\] It must therefore use exactly \(s\) short corridors and project to
the unique \(H\)-geodesic \(P\).

Because \(P\) alternates colours, a switch is necessary between every
two successive corridors. Consequently, \[
d_{G_g}((v_0)_R,(v_s)_R)\ge Ss+s-1.
\qquad\text{(6.11)}
\]

Combining (6.3), (6.10), and (6.11), \[
Ss+s-1-E
\le d_\lambda(v_0,v_s)
\le \lambda(P)
\le Ss+2+E.
\] Therefore \[
s-3\le2E,
\] contradicting (6.5).

This proves that (6.1) is impossible.

\begin{center}\rule{0.5\linewidth}{0.5pt}\end{center}

\subsubsection{7. Conclusion and scope}\label{conclusion-and-scope}

For the fixed class \(\mathcal C\) defined in (1.1), every \(g\ge10000\)
gives an onto \((2,1)\)-quasi-isometry \[
G_g\to H_g^+\in\mathcal C,
\] but every \((1,A)\)-quasi-isometry from \(G_g\) to any member of
\(\mathcal C\) satisfies \[
A+1>\frac g{10000}.
\]

Hence the conjecture, with exactly the closure assumptions in the
question, is \textbf{false for finite graphs}.

The crucial addition to the supplied attempt is Lemma 3.1: it turns an
arbitrary additive quasi-isometry to an arbitrary member of
\(\mathcal C\) into edge lengths on the prescribed complete-graph
skeleton, with only a constant-factor increase in additive error. The
long corridors enforce this rigidity; the red/blue averaging then
supplies the contradiction.

No infinite-graph convention or external conjecture is needed. A version
imposing the additional requirement of closure under edge deletion is
not addressed by this counterexample.

\end{document}
